% TOP_PROBLEM: {"id": 3100086, "title": "Is the exponent of matrix multiplication 2", "category_id": 2, "category": {"id": 2, "name": "combinatorics", "display_name": "Combinatorics", "description": "Counting problems, graph theory, discrete structures.", "slug": "combinatorics", "order_index": 2, "created_at": "2026-07-31T15:26:25.671Z"}, "status": "partially_solved"}
% FIRST_POSTED: 2026-09-27
% !TeX program = lualatex
% Standalone notebook; compile twice: lualatex 3100086.tex
% Original section preserved verbatim except for marked insertion blocks.
\documentclass[11pt,a4paper]{article}
\usepackage[margin=24mm,headheight=14pt]{geometry}
\usepackage{fontspec}
\setmainfont{Latin Modern Roman}
\setsansfont{Latin Modern Sans}
\setmonofont{Latin Modern Mono}
\usepackage{amsmath,amssymb,mathtools}
\usepackage{unicode-math}
\setmathfont{Latin Modern Math}
\usepackage{microtype}
\usepackage{enumitem,xurl,needspace}
\usepackage[dvipsnames]{xcolor}
\usepackage{hyperref}
\hypersetup{unicode=true,colorlinks=true,linkcolor=MidnightBlue,urlcolor=MidnightBlue,
 pdftitle={Research notebook: Problem 3100086},
 pdfauthor={Open Problems Math research notebook},bookmarksnumbered=true}
\usepackage{bookmark,titlesec,fancyhdr}
\titleformat{\section}{\Large\bfseries\raggedright}{\thesection}{0.7em}{}
\pagestyle{fancy}\fancyhf{}
\fancyhead[L]{\small\sffamily Open Problems Math\enspace|\enspace Research notebook}
\fancyhead[R]{\small\sffamily Problem 3100086}
\fancyfoot[L]{\small\sffamily 27 September 2026}
\fancyfoot[R]{\small\sffamily \thepage}
\setlength{\parindent}{0pt}
\setlength{\parskip}{5pt plus 1pt}
\setlength{\emergencystretch}{3em}
\setcounter{secnumdepth}{1}
\setcounter{tocdepth}{1}
\setlist[itemize]{leftmargin=3em,itemsep=5pt,topsep=4pt,parsep=0pt}
\allowdisplaybreaks
\newcommand{\N}{\mathbb{N}}
\newcommand{\Z}{\mathbb{Z}}
\newcommand{\Q}{\mathbb{Q}}
\newcommand{\R}{\mathbb{R}}
\newcommand{\C}{\mathbb{C}}
\newcommand{\F}{\mathbb{F}}
\newcommand{\A}{\mathbb{A}}
\newcommand{\PP}{\mathbb P}
\newcommand{\E}{\mathbb{E}}
\newcommand{\eps}{\varepsilon}
\newcommand{\dd}{\,\mathrm d}
\newcommand{\sref}[2]{\hyperlink{source:#1:#2}{[S#2]}}
\newcommand{\eref}[2]{\hyperlink{extra:#1:#2}{[E#2]}}
\newif\ifshowcatalogue
\showcataloguetrue
\newcommand{\cataloguescope}[1]{\ifshowcatalogue
 \paragraph{Exact scope in the supplied catalogue.}
 \begin{quote}\small #1\end{quote}\fi}
\numberwithin{equation}{section}
\begin{document}
\setcounter{section}{0}
% BEGIN PRESERVED SECTION
% ================================================================
% problemId: 3100086
% Canonical problem ID: 3100086
% exactTarget SHA-256: 83a5afc56933a0edb05d3da6925c3ce585879c630d39b7c712d10ca9ee938d18
\clearpage
\section*{\texorpdfstring{Is the exponent of matrix multiplication 2}{Is the exponent of matrix multiplication 2}}\label{3100086}
\subsection{Definitions and mathematical statement}
The rational bilinear rank $R_{\Q}(n)$ is the least $r$ admitting rational linear forms $\alpha_j,\beta_j$ and fixed rational matrices $C_j$ with
\[
 AB=\sum_{j=1}^r\alpha_j(A)\beta_j(B)C_j\qquad(A,B\in\Q^{n\times n}).
\]
Set $\omega_{\Q}=\inf\{\tau:R_{\Q}(n)=O(n^\tau)\}$. The target is $\omega_{\Q}=2$, or equivalently the existence, for every $\varepsilon>0$, of exact arithmetic circuits of size $O_\varepsilon(n^{2+\varepsilon})$. This counts arithmetic operations at unit cost. It neither bounds coefficient bit lengths nor asserts that the infimum is attained by an $O(n^2)$ algorithm.
\par\smallskip\textit{Formulation and concept references:} \sref{3100086}{1}.
\cataloguescope{Let R\_Q(n) be the tensor rank over Q of the bilinear map multiplying two n by n rational matrices: the least r for which AB is a sum of r products alpha\_j(A) beta\_j(B) C\_j, with rational linear forms alpha\_j, beta\_j and rational output matrices C\_j. Set omega\_Q = inf\{tau in R : R\_Q(n) = O(n\textasciicircum{}tau)\}. Prove omega\_Q = 2. Equivalently, for every epsilon > 0, exact multiplication admits division-free arithmetic circuits over Q of size O(n\textasciicircum{}(2+epsilon)), counting rational additions, subtractions, multiplications and constant scalar operations at unit cost. Circuits need not be uniformly generated and rational coefficient bit lengths are not charged. Constants and thresholds may depend on epsilon. Literal O(n\textasciicircum{}2) attainment is not required.}
\subsection{Short English statement}
Can exact rational matrix multiplication be performed with an operation count arbitrarily close to quadratic?
% BEGIN RESEARCH 26
\subsection{Research attempt}

\paragraph{Current frontier.}
Bl\"aser's survey explains why exact rank, approximation by polynomial
families, and the asymptotic exponent must be separated. \sref{3100086}{1},
Sections~2 and~6. Dupont and collaborators' August 2026 manuscript reports
$\omega<2.371177$ by optimizing combination-loss analysis; its numerical
optimization and full certification have not been independently reproduced
here. \sref{3100086}{2}, \eref{3100086}{1}. Blasiak--Cohn--Grochow--Pratt--Umans
construct a framework for finite algorithms from infinite groups; its
remaining separating-polynomial and dimension requirements are substantive,
not already a proof of exponent two. \sref{3100086}{3}, Introduction.
Alman's universal-method barrier for Coppersmith--Winograd tensors
precludes reaching two within that stated framework, not by every possible
matrix multiplication method. \eref{3100086}{2}.

\paragraph{Attack route.}
The selected route is to combine degenerations with field descent. A
candidate discovered over $\C$ or a number field is useful to this exact
rational problem only after a legitimate conversion. Likewise a limiting
low-rank formula must be converted into exact arithmetic before it bounds
$R_{\Q}$. I make both conversions explicit and ask whether their overhead
could be the missing obstacle. It is not: it can be amortized. The remaining
problem is to find the low-rank family itself, rather than to round
coefficients or set a small parameter equal to zero.

\paragraph{Attempt.}
Let $T_n$ be the rational matrix multiplication tensor. Tensor rank here
counts products of one linear form in each input, with an output vector;
all coefficients must lie in the stated field. First,
\begin{equation}\label{r26:basic}
 R_{\Q}(n)\ge n^2,\qquad T_a\otimes T_b\cong T_{ab},
 \qquad R_{\Q}(ab)\le R_{\Q}(a)R_{\Q}(b).
\end{equation}
For the lower bound, the output matrices in any rank decomposition must
span every matrix $AB$; taking $A=I$ shows their span has dimension at
least $n^2$. For the tensor identity, index a row or column by a pair
of indices and multiply elementary matrices. Expanding two rank
decompositions proves submultiplicativity. These arguments give
$\omega_{\Q}\ge2$, not an upper bound near two.

\textbf{Lemma 26.1 (exact extraction from a rational degeneration).}
Suppose for some fixed $n\ge2$ there is a polynomial tensor identity over
$\Q[t]$
\begin{equation}\label{r26:degeneration}
 D(t)=\sum_{i=1}^{r}a_i(t)\otimes b_i(t)\otimes c_i(t)
          =t^aT_n+\sum_{j=a+1}^{b}t^jU_j,
          \qquad 0\le a\le b,
\end{equation}
where the vector-valued factors are polynomials. Then for every $k\ge1$,
\begin{equation}\label{r26:exact}
       R_{\Q}(n^k)\le(bk+1)r^k,
       \qquad \omega_{\Q}\le\log_n r.
\end{equation}
\emph{Proof.} The coefficient of $t^{ak}$ in $D(t)^{\otimes k}$ is
exactly $T_n^{\otimes k}$: every factor has lowest degree $a$, so there
are no other contributions at the lowest total degree. The whole
polynomial has degree at most $bk$. Choose $bk+1$ distinct rational
numbers $t_0,\ldots,t_{bk}$. Polynomial interpolation expresses the
coefficient of $t^{ak}$ as a rational linear combination of the values
$D(t_l)^{\otimes k}$. Each such value has rank at most $r^k$.
Multiplying one tensor factor by the corresponding rational scalar and
summing proves the first assertion, including the case $b=0$.

For any matrix size $N$, choose $k$ with $n^{k-1}<N\le n^k$, pad with
zeros to size $n^k$, and project back. The bound is
$O(\log(N+1)N^{\log_n r})$, with constants depending on the fixed base
identity. For every $\varepsilon>0$, the logarithm is
$O_\varepsilon(N^\varepsilon)$, so the definition as an infimum gives
the second assertion. No claim that the logarithm disappears in an
exact attained exponent is made.\hfill$\square$

This is a direct version of the approximation-to-exact mechanism in
\sref{3100086}{1}, Section~6, proved here to make the coefficient field and
parameter losses explicit. The degrees $a,b$ may be very large for the
chosen base tensor but are fixed before $k\to\infty$. Allowing them to
grow unaccountably with $k$ would invalidate the estimate.

\textbf{Lemma 26.2 (a finite extension costs only a constant after tensoring).}
If $L/\Q$ has degree $d$ and a rational bilinear tensor $T$ has rank at
most $r$ over $L$, then
\begin{equation}\label{r26:descent}
 R_{\Q}(T)\le d^2r,
 \qquad R_{\Q}(T^{\otimes k})\le d^2r^k.
\end{equation}
\emph{Proof.} Fix a $\Q$-basis $e_1,\ldots,e_d$ of $L$ and a
$\Q$-linear map $\lambda:L\to\Q$ that is the identity on $\Q$.
For rational input vectors, write the two input linear forms in a
rank-one summand as $\alpha=\sum_u e_u\alpha_u$ and
$\beta=\sum_v e_v\beta_v$, with rational linear forms $\alpha_u,\beta_v$.
Applying $\lambda$ coordinatewise to its output gives
\[
          \lambda(C\alpha\beta)
               =\sum_{u,v}\alpha_u\beta_v\,\lambda(Ce_ue_v),
\]
a sum of at most $d^2$ rational rank-one terms. Applying this to all
$r$ summands preserves $T$ because its output on rational inputs is
rational. This proves the first bound. For the second, first tensor the
rank-$r$ decomposition over the \emph{same} field $L$, and only then
apply the first bound. The resulting overhead is $d^2$, not $d^{2k}$.
\hfill$\square$

It follows from the tensor identity and padding that any one exact
rank-$r$ multiplication algorithm of base size $n$ over a number field
gives $\omega_{\Q}\le\log_n r$. Coefficient bit lengths are not being
estimated; the target explicitly does not charge them.

Even a complex exact decomposition of fixed size can be replaced by
one over some number field. Its finitely many coefficient identities
form polynomial equations over $\Q$. A complex solution means their
ideal is proper. The weak Nullstellensatz over $\Q$ gives a maximal
ideal above it whose residue field is a finite algebraic extension of
$\Q$, and hence a solution over a number field. \eref{3100086}{3}.
Equivalently, one can solve over $\overline\Q$ and collect the finitely
many coordinates in one number field. This is an existence argument,
not an efficient coefficient-recovery algorithm.

Thus the arithmetic-operation exponent does not gain an obstacle merely
from passing from $\C$ to $\Q$. More explicitly, if a complex family has
$R_{\C}(N)=O(N^\tau)$, choose a sufficiently large \emph{fixed} base
$n$ with $\log_n R_{\C}(n)\le\tau+\varepsilon$. Descend its one
finite decomposition and tensor it as above; the extension degree is
then fixed. This gives $\omega_{\Q}\le\tau+\varepsilon$. The reverse
inequality follows because a rational decomposition is complex too.
Taking infima establishes equality of these exponents, without claiming
equality of each finite-size rank.

\medskip
\textbf{A concrete benchmark rather than a new upper bound.}
For
$A=\left(\begin{smallmatrix}a&b\\c&d\end{smallmatrix}\right)$ and
$B=\left(\begin{smallmatrix}e&f\\g&h\end{smallmatrix}\right)$, take
\[
\begin{array}{lll}
 M_1=(a+d)(e+h),&M_2=(c+d)e,&M_3=a(f-h),\\
 M_4=d(g-e),&M_5=(a+b)h,&M_6=(c-a)(e+f),\\
 M_7=(b-d)(g+h).&&
\end{array}
\]
The four outputs
\[
 (M_1+M_4-M_5+M_7,\ M_3+M_5,\ M_2+M_4,\ M_1-M_2+M_3+M_6)
\]
are exactly $(ae+bg,af+bh,ce+dg,cf+dh)$. Expanding proves the identity,
and the companion script independently verifies it symbolically. This
is Strassen's known seven-product formula, discussed in \sref{3100086}{1},
not a newly found algorithm. Tensoring yields $\log_2 7>2$; repeated
use of one fixed formula cannot make that exponent drift toward two.

The positive route still needs, for every $\varepsilon>0$, a base size
$n\ge2$ and a certified exact or polynomial-degeneration rank
$r\le n^{2+\varepsilon}$. A bounded improvement for one base size does
not provide this family. Neither interpolation nor field descent reduces
$r$; each only transfers an existing identity into the target model.

\paragraph{Stress test.}
A floating-point approximate identity is not
\eqref{r26:degeneration}. The latter requires the low-degree coefficients
to agree \emph{exactly}. A small nonzero coefficient error may survive
tensor powers and interpolation. Simply sending a numerical parameter
to zero can also make the individual linear forms diverge. Polynomial
coefficient extraction avoids that issue only after a genuine algebraic
identity has been established. The exact interpolation calculation on
a small polynomial tensor is included in the companion checks.

The fixed-field qualification in Lemma~26.2 is essential to its useful
asymptotics. Paying for a new unrelated extension at every tensor level,
or charging enormous coefficient descriptions as bit operations, would
be a different complexity analysis. The statement intentionally follows
the given nonuniform unit-cost rational model. No efficient procedure
for discovering the decompositions is required by that target.

The universal-method barrier in \eref{3100086}{2} shows that merely continuing
to optimize a Coppersmith--Winograd-based implementation within that
method cannot reach two; it does not forbid the infinite-group framework
in every form or a different tensor family. The infinite-group paper
identifies an explicit dimension/separating-degree trade-off still to be
met. \sref{3100086}{3}, Abstract and Section~4. No unproved construction
satisfying that trade-off has been inserted here. The target also would
not, by itself, decide Boolean P versus NP: its unit-cost arithmetic
model and its single bilinear problem are different.

\Needspace{8\baselineskip}
\paragraph{Outcome.}
\textbf{Outcome: Exploratory approach with unresolved gap.}

The notebook gives fully explicit exactness and field-descent conversions,
including tensor-power and all-size overheads, and tests them on an exact
benchmark. These are standard mechanisms developed here in the precise
rational model, not claimed historical novelties. They show that numerical
rounding, coefficient fields, and a failure to attain literal quadratic
rank are not valid substitutes for the real missing step: a sequence of
certified decompositions approaching quadratic exponent. No improved
rank family or proof that such a family cannot exist has been obtained.

% END RESEARCH 26

\subsection{Sources}
\begin{itemize}\raggedright
\item[\textnormal{[S1]}]\hypertarget{source:3100086:1}{} Markus Bläser, Fast Matrix Multiplication, Theory of Computing Graduate Surveys 5 (2013).
\url{https://theoryofcomputing.org/articles/gs005/gs005.pdf}.
{\small\textit{Catalogue formulation source.}}
\item[\textnormal{[S2]}]\hypertarget{source:3100086:2}{} Improving the matrix multiplication exponent with modern optimization and AlphaEvolve.
\url{https://arxiv.org/abs/2608.16884}.
{\small\textit{Catalogue research/status reference; not a proof endorsement.}}
\item[\textnormal{[S3]}]\hypertarget{source:3100086:3}{} Finite Matrix Multiplication Algorithms from Infinite Groups.
\url{https://drops.dagstuhl.de/storage/00lipics/lipics-vol325-itcs2025/html/LIPIcs.ITCS.2025.18/LIPIcs.ITCS.2025.18.html}.
{\small\textit{Catalogue research/status reference; not a proof endorsement.}}
% BEGIN ADDED REFERENCES 26
\item[\textnormal{[E1]}]\hypertarget{extra:3100086:1}{} Emilien Dupont et al.,
\emph{Improving the matrix multiplication exponent with modern optimization
and AlphaEvolve}, arXiv:2608.16884, 17 August 2026.
\url{https://arxiv.org/abs/2608.16884}.
{\small\textit{Version-specific supplement to [S2]; reported bound
$2.371177$, not an independently reproduced optimization certificate.
Consulted 27 September 2026.}}
\item[\textnormal{[E2]}]\hypertarget{extra:3100086:2}{} Josh Alman,
\emph{Limits on the Universal Method for Matrix Multiplication},
arXiv:1812.08731v2, 1 May 2019.
\url{https://arxiv.org/abs/1812.08731v2}.
{\small\textit{Abstract and stated Coppersmith--Winograd universal-method
barrier; no lower bound against all algorithms is inferred.
Consulted 27 September 2026.}}
\item[\textnormal{[E3]}]\hypertarget{extra:3100086:3}{} The Stacks Project,
\emph{Hilbert Nullstellensatz}, Theorem~10.34.1, tag~00FV.
\url{https://stacks.math.columbia.edu/tag/00FV}.
{\small\textit{Finite algebraic residue fields of maximal ideals in a
polynomial algebra; used for fixed finite rank-decomposition equations.
Consulted 27 September 2026.}}

% END ADDED REFERENCES 26
\end{itemize}

% END PRESERVED SECTION
\end{document}
